% =====================================================================
\section{Reproducibility}
\label{app:repro}
% =====================================================================

Every number in this paper that is not proved in the text is produced by a
script in the accompanying supplementary package, and every such number is
attributable to a line of a committed log. This appendix records what the
package contains, how to run it, and which numbers come from where.

\subsection{Package layout}

\begin{verbatim}
PAPER/supplementary/
  README.md                  index of the package
  ENVIRONMENT.md             exact environment used
  MANIFEST.sha256            SHA-256 of every file in the package
  scripts/                   the computation scripts (copies of the
                             committed 04_COMPUTATION/ scripts)
  logs/                      the committed logs the numbers are read from
  generator_matrices/        generator matrices of the headline codes,
                             one .json per code
  parameter_tables/          machine-readable versions of the tables of
                             this paper (CSV)
  CERTIFICATION_PROCEDURE.md what "certified" means here, and how to
                             re-certify a code from scratch
\end{verbatim}

The scripts are byte-identical copies of the scripts committed in
\texttt{04\_COMPUTATION/} of the project repository; the logs are the
committed logs. Nothing was re-run to produce the package, so no script hash
is invalidated.

\subsection{Environment}

The computation is pure Python~3 standard library. No third-party package
(numpy, sympy, GAP, Magma) is used or required; arithmetic in $\F_q$ and in
$A_t$, $R$, $\Rn$ is implemented in \texttt{matrix\_ring\_lib.py} and is
elementary by design, so that every step can be inspected. See
\texttt{ENVIRONMENT.md} for the interpreter version and platform of the
recorded run.

\subsection{Randomness}

The search uses pseudo-random sampling for the families that are too large to
enumerate. The seed is fixed (\texttt{SEED = 20261210}) and every random
choice is drawn from a single \texttt{random.Random(SEED)} instance in a
fixed order of families, so a re-run reproduces every line of the log except
the per-family elapsed-time annotations and the peak-memory line, which are
wall-clock and resource readings and take no part in any computation.

The principal results of this paper do not depend on the sampling: the
theorems are proved, and the headline codes come from families that are
enumerated exhaustively.

\subsection{Where the numbers come from}

\begin{itemize}[leftmargin=*]
\item Theorems~\ref{thm:duality}--\ref{thm:parameters},
  \ref{thm:cosetdistance}, \ref{thm:cosetsingleton},
  \ref{thm:feasibility}, \ref{thm:bookkeeping}, \ref{thm:gram},
  \ref{thm:obstruction}, \ref{thm:hullso}, \ref{thm:fpbasis},
  \ref{thm:nonprincipal}: proved in the text; no computation enters.
\item Table~\ref{tab:witnesses} (the $\MU_{\Bq}=3$ witnesses) and the four
  witness codes of Appendix~\ref{app:tables}: Part~5 of
  \texttt{phase7\_extended\_search\_log.txt}, produced by
  \texttt{phase7\_extended\_search.py}.
\item Table~\ref{tab:cert} (all 28 certified codes): Part~6 of the same log.
\item Table~\ref{tab:headline}: Part~4(d) (Pareto frontier), Part~4(e) (best
  per route) and Part~6 of the same log.
\item Table~\ref{tab:budget} (minimum entanglement per distance): Part~4(c)
  of the same log.
\item Table~\ref{tab:comparison} and Table~\ref{tab:grassl}: Part~8 of the
  same log for the $q=4$ rows, and the pages of \cite{GrasslCodetables}
  retrieved on 2026-11-04 for the $q=9$ rows and for the bounds themselves.
\item The counts $120\,721$, $5\,055$, $4\,418$, $637$, $368$, $88$, $280$,
  $25$ families scanned, $42$ families in the census ($23$ exhaustive, $4$
  counted-only, $15$ blocked): the completion report and Part~2 of the same
  log.
\item The bookkeeping identities of Theorem~\ref{thm:bookkeeping} are
  proved, and check A1 of the log verifies them on all $5\,055$ codes
  ($0$ violations). The coset Singleton bound is check H1 ($0$ violations).
  The Singleton and Hamming bounds are check G1 ($0$ violations).
\end{itemize}

\subsection{How to reproduce}

\begin{verbatim}
  cd PAPER/supplementary/scripts
  python3 phase7_extended_search.py > out.txt 2>&1
  diff <(grep -v '\[[0-9]*s\]' out.txt) \
       <(grep -v '\[[0-9]*s\]' ../logs/phase7_extended_search_log.txt)
\end{verbatim}

The run takes about $9$ minutes and about $370$~MB of resident memory on the
recorded platform. The \texttt{grep -v} filter removes the elapsed-time
annotations; the peak-memory line is the only other difference.

The earlier phases are reproduced the same way from their own scripts; see
\texttt{04\_COMPUTATION/README.md} in the repository for the dependency
order (\texttt{matrix\_ring\_lib.py} first, then the phase scripts).

\subsection{What is \emph{not} reproducible from the package}

The $q=9$ bounds of Table~\ref{tab:grassl} were read from a web page that
changes; the retrieved values are dated 2026-11-04 and are quoted verbatim,
and the page reported ``Last change: 10.06.2024''. They are external data,
not an output of this computation.

\subsection{Verification of the package itself}

\begin{verbatim}
  cd PAPER/supplementary
  sha256sum -c MANIFEST.sha256
\end{verbatim}

Every file listed in the manifest is checked against its recorded hash; see
\texttt{PAPER/SHA256\_MANIFEST.md} for the same check over the manuscript and
the whole package.
